\paragraph{First-order equivariant class: one loop about commuting matrices, the reflected form on labels and the sign of the torsion modes (conditional; bounded no-go).}
The preceding paragraph leaves open which form governs the label modes. Two candidates are tested here at the flat point: a form induced at one loop, and the form that a reflection symmetry of the matrix action assigns to them.

\emph{One loop.} Supply the content of~\cite{IKKT1997}, an input since Equation~\eqref{eq:Smatrix} is bosonic: the ten-component vector with weight \(1\), the ghosts with weight \(-2\) and one sixteen-component Weyl spinor with weight \(-\tfrac12\), with the operators \(H_i=-\Box_{\hat X}\otimes1+\Sigma_i^{ab}\otimes\mathrm{ad}[\hat X_a,\hat X_b]\) of the paragraph on the supertrace at an arbitrary matrix fluctuation in Section~\ref{sec:status-F}, now for operators or finite matrices rather than Poisson brackets, and let \(\Gamma_\varsigma=\tfrac12\sum_iw_i\operatorname{Tr}\ln(H_i+\varsigma)\) with one shift \(\varsigma\) common to all species, in background Feynman gauge off shell, at vanishing fermion background and without the phase of the fermion determinant. With diagonal fermion zero modes, integrating out the off-diagonal fermions gives a nonzero quadratic term in the off-diagonal \(A\)~\cite{AokiIsoKawaiKitazawaTada1998}. The species identities there, \(\sum_iw_id_i=0\) and \(\sum_iw_i\operatorname{tr}(\Sigma_i\cdots\Sigma_i)=0\) with up to three factors, concern constant matrices, and the Duhamel expansion needs no commutation of its operators. For \(\hat X_a=\hat X^{(0)}_a+A_a\) with \([\hat X^{(0)}_a,\hat X^{(0)}_b]=0\) each commutator is of first order in \(A\), so
\[
\Gamma_\varsigma[\hat X^{(0)}+A]-\Gamma_\varsigma[\hat X^{(0)}]=O(A^4)
\]
at finite \(N\) and fixed \(\varsigma\neq0\), for all kernels \(A_a(m,m')\) and all \(Q\), and formally for operators. The cancellation to third order in the commutators is the operator-level statement of~\cite{IKKT1997}, with the first nonvanishing term at fourth order also in~\cite{Steinacker2023}; the equation above is its corollary for commuting backgrounds. At \(N=3\), with random data, the difference scales with exponent \(4.0\) in the amplitude for the Euclidean contraction at \(\varsigma=0.3\) and \(1\), and for the Lorentzian \(\eta\) with complex \(\varsigma\) (numerical); the coefficient depends on \(\varsigma\): at the data tested it grows as \(\varsigma^{-2}\) when \(\varsigma\to0\), through the modes with \(Q=0\), and the exponent \(4.0\) is reached only for amplitudes small compared with \(\varsigma^{1/2}\) (numerical), so the limits \(A\to0\) and \(\varsigma\to0\) do not commute and nothing is stated at \(\varsigma=0\). This trace includes the modes with \(Q=0\), which the standard one-loop object excludes as zero modes~\cite{IKKT1997,AokiIsoKawaiKitazawaTada1998}; no statement is made about that object. So at one loop, with this content and a common regulator at fixed \(\varsigma\neq0\), the second variation about commuting matrices vanishes, and no mass, kinetic form or term of higher order in \(Q\) is induced on the labels at second order in the fluctuation; the form at this order remains the tree form. About commuting matrices this gives at all \(Q\) the vanishing, at this content, of the mass term found in~\cite{Sakai2017} in the semiclassical limit at vanishing base momentum, with coefficient proportional to \(d-2-2^{[d/2]-\kappa}\), which vanishes for \(d=10\), \(\kappa=2\), and which that reference finds positive when bosons exceed fermions. The extension uses the content: with \(\sum_iw_id_i=0\) alone (\(d=4\), one Majorana spinor) the second-order term is nonzero (exponent \(2.0\)), and without fermions there is a term linear in the diagonal part of \(A\), the potential of the eigenvalues, and a second-order term for off-diagonal \(A\). The mass from quantum corrections suggested in~\cite{FurutaHanadaKawaiKimura2007} therefore does not arise at this order about commuting matrices with this content and regulator. A second-order term requires nonzero commutators of the background: for four commuting matrices tensored with \(\mathbb 1\) on \(\mathbb C^2\), three further matrices \(\mathbb 1\otimes\sigma^i/2\) and \(A\) along the first four, the exponent is \(2.0\) at this content. There the second-order terms are those of the factorisation in the paragraph on product backgrounds, and their value on the labels is not computed.

\emph{Commutant labels at nonzero torsion.} If \(\mathcal C\) commutes with all \(\hat X_a\) on any background, the adjoint compatibility of the preceding paragraph gives \(\mathcal E_b[\tfrac12\{f,\mathcal C\}]=0\) for constant \(f\), since \(\mathcal E_b\) is built from \(\mathrm{ad}\,\hat X_a\) and \(\mathrm{ad}[\hat X_a,\hat X_b]\); a commutant label therefore acquires no term without a derivative of \(f\) in its linearised equation, background torsion included. Such a term can come only from a mode that stops being commutant. On \(\hat X_a=\mathrm iu_a{}^bE_b\), with \(E_b\) left-invariant fields on a four-dimensional Lie group, structure constants \(\mathsf c_{ab}{}^d\) and adjoint matrices \(\mathrm{ad}_a\), one has \(\mathsf R_{ab}{}^{cd}=0\) and \(\mathsf T\) built from \(\mathsf c\), and the matrix equation becomes \(\eta^{ac}\mathrm{ad}_a\mathrm{ad}_c=0\), which has the form of the equation \(h_{ab}[t^a,[t^b,t^c]]=0\) of~\cite{Liao2026} for a form \(h\) on the algebra, with the indices placed oppositely there; Hermiticity constants change no commutator. An operator commuting with all \(\hat X_a\) that is a right translation of \(u\) by \(\lambda\) plus a vector field \(V_\lambda\) on the group needs \([E_b,V_\lambda]=\lambda_b{}^cE_c\), so \(\lambda\) must be a derivation of the algebra. For \(E_0-E_3\) central, \([E_0+E_3,E_1]=E_2\) and \([E_0+E_3,E_2]=-E_1\), these \(\lambda\) in \(\mathfrak{so}(3,1)\) form a line, and for \([E_0,E_2]=E_2\), \([E_0,E_3]=E_3\), \([E_1,E_2]=E_2\), \([E_1,E_3]=-E_3\) only \(\lambda=0\) (exact), against six at \(\mathsf c=0\). Whether the remaining modes acquire a mass term of definite sign is not computed.

\emph{Reflected form.} Use the Euclidean contraction and the compact fibre \(\mathrm{SO}(4)\), as for \(\mathrm{Spin}(n)\) in~\cite{Hanada2006}, so that \(\hat X_a=-u_a{}^bp_b\) on \(L^2(\mathbb R^4\times\mathrm{SO}(4))\). Let \(\theta=\mathrm{diag}(-1,1,1,1)\), with entries \(\theta_r\), \(U_\theta\psi(p,u)=\psi(\theta p,\theta u\theta)\) and \(\mathcal I[A]_a=\theta_a{}^cU_\theta A_cU_\theta^{-1}\). Then \(U_\theta\hat X_aU_\theta^{-1}=\theta_a{}^c\hat X_c\), and \(\mathcal I\) is a unitary involution of the Hilbert--Schmidt kernels which preserves the background and the action of Equation~\eqref{eq:Smatrix} at \(C_{\mu\nu\lambda}=0\). The form \(\mathrm{HS}(A,\mathcal I[B])\) is therefore indefinite and nondegenerate, with the Hilbert--Schmidt form as its majorant in the sense of Krein spaces. Let \(w\ge0\) be a function of the mean momentum, invariant under \(\theta\), of unit total mass and finite second moments, and define for \(r,g=0,\dots,3\), with the kernel \(A_a\) between \(m=(\bar p+q/2,u)\) and \(m'=(\bar p-q/2,u)\),
\[
\phi_r{}^g(y)=\int\mathrm d\mathrm{Haar}(u)\int\mathrm d^4\bar p\,\frac{\mathrm d^4q}{(2\pi)^4}\;e^{\mathrm iq\cdot y}\,\bar p_g\,w(\bar p)^{1/2}\,u^a{}_rA_a.
\]
A gauge variation \(A_a=Q_a\Lambda\) shifts \(\phi_r{}^g\) by a multiple of the gradient \(\partial_r\Lambda^g\), with \(\Lambda^g\) the same projection of \(\Lambda\), so \(\partial_r\phi_s{}^g-\partial_s\phi_r{}^g\) is invariant at this order, and \(\phi_r{}^g\) evaluated on \(\mathcal I[A]\) is \(\theta_r\theta_g\,\phi_r{}^g(\theta y)\) evaluated on \(A\). The Gaussian integral of the background-gauge form gives, for \(y\neq y'\),
\[
\bigl\langle\phi_r{}^g(y)\,\phi_s{}^h(y')\bigr\rangle=g_{\mathrm{YM}}^2\,\mathcal V\,\delta_{rs}\,\mathsf G_{gh}\,\frac1{4\pi^2|y-y'|^2},
\]
with \(\mathsf G_{gh}=\int w\,p_gp_h\) and \(\mathcal V\) a positive factor common to all indices that contains the fibre delta at \(u=u'\), so it is singular; smearing the fibre to make it finite changes \(Q\), and that case is not treated. For a test function supported at \(y^0>0\) the reflected form of \(\phi_r{}^g\) is therefore \(\theta_r\theta_g\mathsf G_{gg}\) times the reflection-positive form of the massless scalar, and that of the invariant field strength with spatial \(r,s\) is \(\theta_g\mathsf G_{gg}\) times the reflection-positive form of the Maxwell field. So on these gauge-invariant probes with \(g=0\) and \(\mathsf G_{00}>0\) the form is negative for generic test functions: the Euclidean flat point with the reflection \(\mathcal I\) does not satisfy the positivity of~\cite{OsterwalderSchrader1973,OsterwalderSchrader1975} on them at tree level. For a weight invariant under \(\mathrm O(4)\), \(\mathsf G_{gh}=c\,\delta_{gh}\) with \(c>0\), and the reflected form on the translations is \(c\,\theta_g\delta_{gh}\), which is \(c\,\eta_{gh}\). For the probes \(\int\mathrm d\Lambda\,\psi_\lambda(\Lambda)R_\Lambda\,w^{1/2}\) of the fibre directions, with \(R_\Lambda\) the right translation of \(u\), \(\psi_\lambda=\tfrac12\lambda^{cd}(\ln\Lambda)_{cd}\,\rho\) and \(\rho\) supported near the identity and invariant under conjugation by \(\mathrm O(4)\), the Gram matrix is a positive multiple of \(\tfrac12\lambda_{cd}\lambda'^{cd}\), the only invariant, and \(\mathcal I\) acts by \(\lambda\mapsto\theta\lambda\theta\), so the reflected pairing of these kernels is the Lorentz invariant form on \(\mathfrak{so}(3,1)\) with boosts negative. A two-point sign is not stated for them at finite smearing, since \(Q\) then depends on \(p\).

The same relation holds for the real Lorentzian fibre of the preceding paragraph. For a Gaussian weight of variance \(P^2\) per component in the rest frame of a unit timelike \(n\), \(\mathsf G_{gh}=P^2(\eta_{gh}+2n_gn_h)\), which is positive and depends on \(n\), and with \(\theta_n\) the reflection along \(n\) one has \(\mathsf G_{gh}(\theta_n)^h{}_k=P^2\eta_{gk}\), independent of \(n\) (exact, also after a boost). So for this weight the positive form of the preceding paragraph is a majorant of the covariant form, the two being related by the involution \(U_{\theta_n}\) on kernels; moments not of this form give no multiple of \(\eta\). This locates the tension recorded in the preceding paragraph. The components with an odd number of time indices, which~\cite{FurutaHanadaKawaiKimura2007} calls negative modes although no action is used there, and does not expect to have positive norm were the equation derived from a field theory action, are those odd under \(\mathcal I\); here their sign is computed for the reflected form at tree level with the matrix action, and it is not a norm computed in a Lorentzian theory. In~\cite{HoKawaiSteinacker2025} the positivity for the planar sector with quenched diagonal elements is read from the large-\(N\) gauge theory with the colour factor removed, the colour degeneracy corresponding to the relative coordinate. The reading adopted here, which that reference does not state in these terms, is that this coordinate, conjugate to \(\bar p\), is inert under the reflection. The corresponding map, \(\theta_a\) with \((Q,\bar p)\mapsto(\theta Q,\bar p)\), preserves the form \(Q\cdot Q|A|^2-|Q\cdot A|^2\) and does not preserve the cubic and quartic terms at finite \(N\) (one finite model, exact matrices), so it is a symmetry of the Gaussian form and not of the matrix action. That reference records no clear correspondence between the bilocal and the tensor inner products; the reflected form supplies one for these truncations, and it does not settle which form is the physical inner product. The no-ghost theorem of~\cite{Steinacker2019NoGhost}, on a different background, gives a positive-definite inner product on gauge-fixed on-shell modes modulo pure gauge; it does not conflict with a tree-level sign on probes at this background.

\emph{Bearing on Section~\ref{sec:status-E}, and the relaxations.} The statement above bounds one route to Lorentzian signature, reconstruction from a Euclidean integral. That route is not asserted in that section, and it is not the route of a Lorentzian matrix integral with a convergence regulator, which is then no longer related to the Euclidean model by Wick rotation~\cite{Steinacker2026Invitation}; the reflection statement does not bear on that integral. Its hypotheses are the Euclidean contraction, a commuting background with a fibre, tree level, reflection by \(\mathcal I\) and probes of the translation labels. Relaxing them gives the following. (i) \emph{Singlets.} For uniform diagonal elements, quenching and planar diagrams the matrix model is the dimensionally reduced gauge theory~\cite{HoKawaiSteinacker2025}. If its singlet correlators in four dimensions satisfy the positivity, which is taken from field theory and not checked here, a Lorentzian signature there is an output of reconstruction and not a chosen rotation. That sector has no frame, and that reference places gravity in the nonplanar diagrams and the diagonal elements, so Assumption~\ref{ass:lorentzian} is not discharged. (ii) \emph{Constraint.} Take \(\eta\) on \(f_b{}^g\) and \(\tfrac12\lambda_{cd}\lambda^{cd}\) on \(f_b{}^{cd}\), at null \(q\neq0\) and for the two polarisations \(b\) transverse to \(q\). On the twenty components these forms have twelve positive and eight negative directions, the eight being those with an odd number of time indices. On \(q_gf_b{}^g=0\) and \(q_cf_b{}^{cd}=0\), the second being the kernel of the rank-three map behind the double zeros, they are nonnegative. With \(v_1,v_2\) transverse unit vectors, the kernels are spanned by \(q,v_1,v_2\) and by \(q\wedge v_1,q\wedge v_2,v_1\wedge v_2\), which gives signature \((2,0)\) with one null direction and \((1,0)\) with two null directions, so six positive and six null directions in all (also exact at three sampled \(q\)). The translation part is the constraint imposed in~\cite{HoKawaiSteinacker2025} without a fibre. Its bivector analogue, for gauge parameters as well, leaves three of six parameters at each momentum and does not close: the coordinate \(y^2\) times \(\mathrm dy^0\wedge\mathrm dy^1\) and the coordinate \(y^3\) times \(\mathrm dy^1\wedge\mathrm dy^2\) have vanishing divergence and their commutator does not. It therefore conflicts with six-parameter local rotations, and it is not derived. (iii) \emph{Mass.} The one-loop result above excludes this at second order about commuting matrices with this content; nonzero background commutators, higher loops and the dynamics of the diagonal elements remain. (iv) \emph{No selection by classical solutions.} The contraction is not selected by the equation at \(\mathsf R_{ab}{}^{cd}=0\): the second algebra above solves it for \(\eta\); \([E_0,E_2]=E_2\), \([E_0,E_3]=E_3\), \([E_1,E_2]=E_3\), \([E_1,E_3]=-E_2\) solves it for \(\delta\) and not for \(\eta\); and the frame \(\partial_t\), \((\alpha-t)\partial_i\) of~\cite{IsonoTomino2010} solves it for both (exact). A classification of Lie-algebraic solutions by signature is in~\cite{Liao2026}.

These statements are for Equation~\eqref{eq:Smatrix} with \(C_{\mu\nu\lambda}=0\) and no matter term, the commuting and left-invariant backgrounds named, formal adjoints with domains and infrared volumes untreated, one common regulator, and tree level except for the first part, where the content is supplied. They do not derive that the matrix integral selects this class, its frame or its non-degeneracy, the Euclidean or the Lorentzian contraction, the weight \(w\) or the quenching, so Assumptions~\ref{ass:frame-lock}, \ref{ass:nondeg}, \ref{ass:lorentzian} and~\ref{ass:local-lorentz} are not discharged. The modes with odd labels have simple zeros at fixed label and for a Gaussian weight the form is diagonal in the label, so the sign above concerns the field content and is not a statement on an Ostrogradsky ghost in the graviton two-point function of Section~\ref{sec:status-N}; Assumption~\ref{ass:EFT-a2} is not discharged. Terms of fourth and higher order, loops beyond the first, a value or sign of \(G_{\mathrm{eff}}\) and the identification of the frame with the emergent metric of Chapter~\ref{ch:vierbein-emergence} are not treated. Problems~D, E, F and~N remain open and PROGRAMME-CONDITIONAL.

